\section{Conclusions}\label{Conclusions}
This paper has outlined a new software package that implements
Bayesian model averaging and jointness measures for \pkg{gretl}. Bayesian
model averaging is a straightforward and natural extension of standard
Bayesian analysis and it is a useful and popular alternative to other
variable selection procedures, especially for a large set of
regressors. Here we used \pkg{gretl}, which is a free, open-source
software for econometric analysis with an easy-to-use GUI. Our goal
was to familiarize potential users with the features and the different
options that our package has to offer. We described how our package
implements Bayesian model averaging, as well as the outputs that are
returned.

Recent trends in Bayesian model averaging still concern applications
in various areas of empirical research, i.e., economic growth (see
\citealt{RePEc:taf:regstd:v:48:y:2014:i:1:p:44-67},
\citealt{GazdaPuziak2012}), stock returns (see
\citealt{RePEc:cqe:wpaper:2914}), geography (see
\citealt{RePEc:eee:csdana:v:75:y:2014:i:c:p:227-238}) and many more.
One can also encounter new papers about theory of Bayesian model
averaging, for example: two-stage Bayesian model averaging (see
\citealt{RePEc:taf:emetrv:v:33:y:2014:i:1-4:p:122-151}) or comparison
of Bayesian model averaging to the weighted average least squares
method (see \citealt{RePEc:dgr:kubcen:2011082}).

The future development of the \pkg{BMA} package could be focused on two
main aspects: speed improvements and feature extensions. The first one
is strongly correlated with the \pkg{gretl} development process, especially
with work on improvements in speed of the \proglang{Hansl} interpreter. It is
possible that the interpreter will be rewritten in the near future to
achieve speed boost with backward compatibility.  The feature
extensions would concern two aspects: implementing new g-priors and
model priors.  In the literature we can find many types of g-priors
(see \citealt{FernandezLeySteel2001a}).  Moreover new model priors
could be also considered.  For instance,
\cite{RePEc:aea:aecrev:v:94:y:2004:i:4:p:813-835} proposed the
so-called fixed model prior.  These implementations would surely
improve the flexibility of the \pkg{BMA} package.
 
